% latexmk -cd -lualatex works/A_11/appendix.tex
% latexmk -c -cd works/A_11/appendix.tex

\documentclass{ees}

\newlist{lyricslist}{description}{1}
\setlist[lyricslist]{
  partopsep=0pt,
  labelindent=0em,
  labelwidth=4.5em,
  leftmargin=4.5em,
  labelsep=0pt,
  first=\ltseries,
  font=\normalfont\itshape\ltseries,
  style=multiline
}

\newenvironment{lyrics}[1]{%
  \subsection{#1}\nopagebreak%
  \begin{lyricslist}%
  \let\voice\item%
}{%
  \end{lyricslist}%
}

\begin{document}

\clearscrheadfoot

\chapter{Lyrics}

\textit{Roles:} Die Barmherzigkeit (S) · Der verlohrene Sohn (A)

\section{Actus primus}

\begin{lyrics}{Scena prima}
  \voice[Der verlohrene Sohn]
  Wo wend ich mich nun hin!\\
  die Noth hat mich umbgeben,\\
  ich bin von aller Weld\\
  verlaßen, und veracht,\\
  weil durch verkherten Sünn\\
  in meinen Wolluſtleben\\
  mein Haabſchafft, Guet und Geld\\
  nun gänzlich durchgebracht,\\
  daß Kleid der erſten Unſchuld\\
  iſt völlig ſchon zerrißen,\\
  mein Erbtheil iſt dahin,\\
  ô Sünder der ich bin!\\
  der Hunger meiner Seelen,\\
  ſo mir ins Gwißen nagt,\\
  will mich zu Todte quällen,\\
  und mir die Speiß verſagt,\\
  ſo doch mein liebſter Vatter\\
  ſein gringſten Knechten gönnet,\\
  ach daß ich doch nur bald\\
  mit ihm wär außgeſöhnet.\\
  Doch nein es iſt zu ſpatt,\\
  mein Ungehorſamb zeiget,\\
  daß ſich zu keiner Gnad\\
  deß Vatters Hertz mehr neiget.
\end{lyrics}

\begin{lyrics}{Scena secunda}
  \voice[Die Barmherzigkeit]
  Nein nein verzage nicht,\\
  du würſt noch Gnade fünden,\\
  dan deines Vatters Hertz\\
  hat allzu bittren Schmertz,\\
  daß du von ihm gewichen,\\
  bekhenne deine Schuld,\\
  werff dich zu ſeinen Füeßen,\\
  rueff nur umb Gnad und Huld,\\
  laß bittre Thränen flüeßen,\\
  ſo würſtu Gnad erlangen,\\
  Verzeühung und Pardon,\\
  Er würd dich ſelbſt umbfangen\\
  als ſeinen liebſten Sohn.
\end{lyrics}

\begin{lyrics}{Aria prima}
  \voice[Barmherzigkeit]
  Ô Menſch dich nicht betrüeb!\\
  dein Gott iſt ſo gearthet,\\
  daß Er nur auff dich warthet,\\
  biß du mit Reu und Schmertz\\
  bewegſt ſein mildes Hertz,\\
  dan Er iſt voll der Lieb.\\[1ex]
  Nun iſt die Gnadenzeit,\\
  wo du Verzeühung fündeſt,\\
  und Gott mit Liebe bündeſt,\\
  daß Er nicht ſtraffen kan\\
  wie Er ſchon offt gethan,\\
  drum dich zur Bueß bereüth!
\end{lyrics}

% \begin{lyrics}{}
%   \voice[]
% \end{lyrics}

% \begin{lyrics}{}
%   \voice[]
% \end{lyrics}

% \begin{lyrics}{}
%   \voice[]
% \end{lyrics}

% \begin{lyrics}{}
%   \voice[]
% \end{lyrics}

% \begin{lyrics}{}
%   \voice[]
% \end{lyrics}

% \begin{lyrics}{}
%   \voice[]
% \end{lyrics}

% \begin{lyrics}{}
%   \voice[]
% \end{lyrics}

% \begin{lyrics}{}
%   \voice[]
% \end{lyrics}

% \begin{lyrics}{}
%   \voice[]
% \end{lyrics}

% \begin{lyrics}{}
%   \voice[]
% \end{lyrics}

% \begin{lyrics}{}
%   \voice[]
% \end{lyrics}

% \begin{lyrics}{}
%   \voice[]
% \end{lyrics}

% \begin{lyrics}{}
%   \voice[]
% \end{lyrics}

% \begin{lyrics}{}
%   \voice[]
% \end{lyrics}

% \begin{lyrics}{}
%   \voice[]
% \end{lyrics}

% \begin{lyrics}{}
%   \voice[]
% \end{lyrics}

% \begin{lyrics}{}
%   \voice[]
% \end{lyrics}

% \begin{lyrics}{}
%   \voice[]
% \end{lyrics}

% \begin{lyrics}{}
%   \voice[]
% \end{lyrics}

% \begin{lyrics}{}
%   \voice[]
% \end{lyrics}

% \begin{lyrics}{}
%   \voice[]
% \end{lyrics}

\end{document}
